% Generated draft: requires mathematical and source review.
\documentclass[11pt]{article}
\usepackage[margin=25mm]{geometry}
\usepackage{fontspec}
\setmainfont{Latin Modern Roman}
\usepackage{amsmath,amssymb,mathtools,unicode-math}
\setmathfont{Latin Modern Math}
\usepackage{enumitem,needspace,xurl,hyperref}
\setlength{\parindent}{0pt}
\setlength{\parskip}{5pt}
\setlength{\emergencystretch}{3em}
\newcommand{\N}{\mathbb{N}}
\newcommand{\Z}{\mathbb{Z}}
\newcommand{\Q}{\mathbb{Q}}
\newcommand{\R}{\mathbb{R}}
\newcommand{\C}{\mathbb{C}}
\newcommand{\F}{\mathbb{F}}
\newcommand{\A}{\mathbb{A}}
\newcommand{\PP}{\mathbb{P}}
\newcommand{\E}{\mathbb{E}}
\newcommand{\eps}{\varepsilon}
\newcommand{\dd}{\,\mathrm{d}}
\begin{document}
\section*{UnsolvedMath ID: 800099}
\textbf{Source problem code:} catalogue entry
\subsection{Definitions and mathematical statement}
Let \(S^1=\R/\Z\), and write \(G=\mathrm{Diff}(S^1)\) for the diffeomorphism group intended in the source. The source does not specify the differentiability class or whether only orientation-preserving diffeomorphisms are allowed; these conventions remain unresolved. Let \(G^\delta\) denote this group with the discrete topology, and let \(BG^\delta\) be its classifying space.

For an explicit integral interpretation, use the bar chain complex \(C_*(G;\Z)\) with trivial coefficients. Its degree-two generators are written \([u\mid v]\), and
\[
\partial [u\mid v]=[v]-[uv]+[u].
\]
Thus, for \(f,g\in G\) satisfying \(fg=gf\), the chain
\[
z(f,g)=[f\mid g]-[g\mid f]
\]
is a cycle. It is a boundary precisely when there exists \(b\in C_3(G;\Z)\) such that
\[
\partial b=z(f,g),
\]
or equivalently when
\[
[z(f,g)]=0\quad\text{in }H_2(BG^\delta;\Z).
\]
\textbf{Question (source Question 40).} For commuting circle diffeomorphisms \(f\) and \(g\), to what extent can one establish that \(z(f,g)\) is a boundary? Determine conditions on the commuting pair under which this vanishing holds.

The integral bar-complex formulation above makes one conventional interpretation explicit; the source does not specify homology coefficients. Its wording also leaves open whether the intended goal is universal vanishing, a characterization of the pairs giving boundaries, or sufficient conditions. These qualifications must be resolved before assigning a more specific conjecture.
\subsection{Short English statement}
For two commuting circle diffeomorphisms, the ordered pairs \((f,g)\) and \((g,f)\) determine a difference that is a two-cycle in the classifying space of the discrete diffeomorphism group. The question asks when this cycle bounds. The regularity, orientation convention, coefficients, and intended scope of the answer are not specified in the supplied source.
\subsection{Sources}
\begin{itemize}
\item Local UnsolvedMath registry, permanent ID \(800099\), exact source code \texttt{catalogue entry}, title ``Foliations and Diffeomorphism Groups --- Question 40.'' The supplied record attributes the problem to Nariman and Yazdi, with proposed year \(2025\), and labels it open; this label is supplied metadata, not an independently verified mathematical assessment.
\item Nariman and Yazdi, \emph{Problem List on Foliations and Diffeomorphism Groups} (2025), Question 40. Supplied source URL: \url{https://arxiv.org/abs/2503.18274}. The record reports access on \texttt{2026-10-16} and extraction by \texttt{arxiv\_latex}; these are supplied provenance details, not a claim of independent source consultation.
\item The supplied question mentions work of Navas, Eynard-Bontemps, and others on commuting circle diffeomorphisms, and describes the question as being posed again forty years after the speaker's paper. No specific references for those works or that earlier paper, and no identification of the first-person speaker, are supplied.
\end{itemize}
\end{document}
